\documentclass{article}

\usepackage[margin=1in]{geometry}
\usepackage{times}
\input{math_commands.tex}

\usepackage{booktabs}
\usepackage{microtype}
\usepackage{natbib}
\usepackage{hyperref}

\title{VoxReason: Source-Cited Speaking Plans for Evidence-Grounded Reasoning}

\author{Anonymous Authors}

\newcommand{\method}{VoxReasonBench}

\begin{document}

\maketitle

\begin{abstract}
Speech models increasingly expose intermediate reasoning before generating or evaluating spoken responses. Yet this reasoning is often judged only through the final answer, leaving open whether the model used the source evidence that governs a speaking plan. We introduce \method, a listener-free source-label diagnostic suite for evidence-grounded speech reasoning. A model must cite permitted source cues, predict a structured speaking plan, and remain stable under counterfactual cue edits. The current release includes source-label cases, public-safe automatic metrics, and reproducible code for evidence precision/recall, decisive-cue recall, plan-slot accuracy, hallucinated-evidence rate, uncited-evidence rate, grounded score, counterfactual consistency, and acoustic-anchor checks. On 100 source-label cases, upper-bound reconstruction exposes a text-only gap of 14.3 points in evidence F1, 100.0 points in decisive-cue recall, and 81.5 points in plan-slot accuracy. Multi-seed model checks further show where larger or process-reward-tuned planners preserve evidence citation while losing plan coverage. \method is not a broad speech benchmark and does not claim listener judgments; instead, it isolates the process-supervision layer that requires verification before human listening studies are meaningful.
\end{abstract}

\section{Introduction}
Speech is not only a transcript. Pitch, timing, energy, speaker role, and conversational state all shape how an utterance is interpreted or spoken. Recent work has begun to place reasoning modules around speech generation and audio understanding, including actor-style speech role-playing, chain-of-thought text-to-speech, deep audio reasoning, audio generation with reasoning traces, and multimodal speech representations~\citep{actormind2026,cottts2026,mmar2025,thinksound2025,prismaudio2026,audiox2026,speechrepresentation2026}. The resulting systems create a new evaluation problem: a fluent explanation may be disconnected from the evidence that justifies the speaking plan.

This paper studies \emph{evidence-grounded speech reasoning}. Given dialogue context, optional acoustic observations, a speaker role, and a target utterance, the system predicts a structured speaking plan and cites the cues that support it. The core question is not whether the final waveform sounds good, but whether the reasoning process identifies the right evidence and changes predictably when that evidence changes.

\method operationalizes this question with a narrow claim boundary. The current paper uses existing source labels and automatic checks; it does not depend on newly recruited listeners. This makes the evaluation practical for rapid method development while avoiding claims that require perceptual human studies.

Our contributions are:
\begin{itemize}
    \item a formal process-supervision target for speech reasoning, separating evidence selection, speaking-plan prediction, and counterfactual cue consistency;
    \item a listener-free evaluation protocol with source-label provenance, leakage controls, and automatic metrics tied to concrete claim boundaries;
    \item a reproducible public code release that rebuilds the main numeric results without site-specific job launch files; and
    \item empirical diagnostics showing that source-label reconstruction exposes a text-only gap while model-scale and process-reward variants reveal remaining plan-coverage weaknesses.
\end{itemize}

\section{Related Work}
\paragraph{Reasoning for speech and audio.}
Actor-style speech role-playing and context-aware speech synthesis motivate explicit intermediate reasoning for expressive speech~\citep{actormind2026,cottts2026}. Broader audio reasoning benchmarks and generation systems examine multimodal chain-of-thought behavior across speech, music, and environmental audio~\citep{mmar2025,thinksound2025,prismaudio2026}. \method differs by making the intermediate process itself the supervised object: each step must cite recoverable cues and survive controlled cue edits. The missing interface is the pre-synthesis decision surface, where typed delivery fields stay linked to source records before any waveform or answer-level score hides the evidence path.

\paragraph{Representation and generation controls.}
Recent audio language modeling work shows that representation choice changes the evidence available for model reasoning~\citep{speechrepresentation2026,codecdoesmatter2025,llaseg12025}. Text-to-audio systems likewise rely on conditioning and representation choices that may confound downstream claims~\citep{audiox2026,spatialaudio2025,sparktts2025}. \method therefore treats representation and source-label checks as part of the evaluation design, not as implementation details.

\section{Task Definition}
Let $x_i=(c_i, a_i, r_i, u_i)$ be one speech-reasoning instance, where $c_i$ is dialogue context, $a_i$ is optional speech evidence, $r_i$ is speaker or role context, and $u_i$ is the target utterance. The supervision contains an evidence set $E_i$ and a structured speaking plan $y_i$ with slots such as emotion, intensity, pitch, energy, rate, pauses, emphasis, and stance. A system predicts cited evidence $\hat{E}_i$ and a plan $\hat{y}_i$.

We evaluate five listener-free quantities. Evidence F1 compares $\hat{E}_i$ with source-labeled cues. Plan-slot accuracy averages slot-level agreement between $\hat{y}_i$ and $y_i$. Hallucinated-evidence rate penalizes citations that are not present in the source labels. Uncited-evidence rate reports required source cues omitted by the prediction. Counterfactual consistency measures whether editing a cue changes only the plan slots licensed by that cue. These metrics target reasoning faithfulness rather than perceived audio quality.

Every result follows a fixed admission rule. A row enters the evidence column only with complete expected predictions, strict schema validity, public-safe provenance, and a listener-free endpoint declared before inspection. Parser checks, lookup ceilings, acoustic anchors, and incomplete learned runs are diagnostic rows. Results that need final seed coverage or listener ratings remain blocked or out of scope.

\section{Method}
\method uses a structured planner that emits machine-checkable JSON. The planner is trained or prompted to cite cue identifiers before predicting plan slots. The same schema supports supervised fine-tuning and process rewards: evidence recall rewards missing-cue recovery, evidence precision penalizes unsupported citations, and slot accuracy rewards correct speaking-plan choices. Counterfactual cue edits provide a stress test: the method is expected to revise affected slots while preserving irrelevant ones. The falsifiable claim is narrow: a planner changes delivery fields only when cited source cues license those changes.

Two design choices are important for reproducibility. First, source labels are kept separate from generated explanations so that invented evidence earns no credit. Second, evaluation is defined before any listener study. This prevents automatic improvements from being presented as user-facing perceptual gains.

\section{Experiments}
\paragraph{Data.}
The public release includes 100 checked source-label cases derived from RAVDESS records~\citep{ravdess2018} for the main diagnostic and 300 acoustic-feature rows for a dependency-free preflight. The labels provide the supervision needed to evaluate evidence and plan slots without recruiting listeners.

\paragraph{Baselines and variants.}
We compare a text-only control against the evidence-grounded planner on source-label cases. We also include three-seed model checks for Qwen2.5-3B SFT, Qwen2.5-7B SFT, and Qwen2.5-7B preference tuning. These checks are diagnostic: they test whether the structured schema and evidence metrics remain stable across model variants.

\paragraph{Metrics.}
All reported metrics are automatic and listener-free. Higher is better for evidence F1, plan-slot accuracy, and grounded score. Lower is better for hallucinated- and uncited-evidence rates. We also compute lightweight acoustic statistics to verify that the speech-side input channel is present and well-formed.

\begin{table}[t]
\centering
\begin{tabular}{lrrr}
\toprule
Metric & Text control & Evidence planner & $\Delta$ \\
\midrule
Evidence F1 & 0.857 & 1.000 & 0.143 \\
Plan acc. & 0.185 & 1.000 & 0.815 \\
Grounded & 0.569 & 1.000 & 0.431 \\
Halluc. rate & 0.000 & 0.000 & 0.000 \\
\bottomrule
\end{tabular}

\caption{Source-label diagnostic on 100 paired cases. Upper-bound reconstruction exposes the text-only gap on process metrics.}
\label{tab:source-label}
\end{table}

\begin{table}[t]
\centering
\begin{tabular}{lrrrr}
\toprule
Model & Evidence F1 & Plan acc. & Grounded & Halluc. rate \\
\midrule
Qwen2.5-3B SFT & 1.000 & 0.811 & 0.915 & 0.000 \\
Qwen2.5-7B SFT & 1.000 & 0.725 & 0.876 & 0.000 \\
Qwen2.5-7B preference & 1.000 & 0.689 & 0.860 & 0.000 \\
\bottomrule
\end{tabular}

\caption{Three-seed listener-free model checks. Evidence citation remains exact in this subset, while plan-slot accuracy reveals the main remaining modeling weakness.}
\label{tab:model-checks}
\end{table}

\begin{table}[t]
\centering
\begin{tabular}{lr}
\toprule
Statistic & Value \\
\midrule
Cases & 300 \\
Feature rows & 300 \\
Duration (s) & 3.770 \\
Silence fraction & 0.570 \\
Voiced fraction & 0.349 \\
Pitch proxy (Hz) & 206.962 \\
\bottomrule
\end{tabular}

\caption{Dependency-free acoustic preflight over 300 public-safe feature rows. These checks verify the speech evidence channel but do not support listener judgments.}
\label{tab:acoustic-preflight}
\end{table}

\section{Results and Analysis}
Table~\ref{tab:source-label} shows that upper-bound reconstruction reaches perfect source-label agreement on the diagnostic cases, while the text-only control misses many plan slots and omits required source cues. The source-label row also reaches 1.000 decisive-cue recall, while the text-only control reaches 0.000 because it does not cite the cue changed by the local counterfactual edit. The largest plan gap is still plan accuracy: reconstruction improves from 0.185 to 1.000. Grounded score improves by 0.431, indicating that the diagnostic gap is not only a citation-format effect but also changes the structured plan.

A construct-validity check gives the boundary of this diagnostic. The public source-label split has zero public context-audio rows, two target utterances, two role labels, one scene label, and fifteen emotion-intensity keys. All fifteen observed keys map to a single gold plan. A deterministic train-set lookup keyed only by source emotion and intensity reaches 1.000 exact-plan accuracy on the test split, while a leave-key-out stress drops to 0.000 exact-plan accuracy and 0.242 plan-slot accuracy. The public code also builds a source-key-disjoint rerun split with 60 train, 16 development, and 24 test cases and zero train/dev/test key overlap. The current release therefore tests source-record accounting and verifier behavior, while the disjoint split is a falsification target for future learned-planner runs rather than completed learned-model ordering evidence.

Table~\ref{tab:model-checks} gives a less flattering but useful view of learned variants. All three variants maintain exact evidence precision/recall on completed predictions, but they leave many examples without valid predictions. The rows are diagnostic rather than ranked: they show that evidence citation may look solved while the plan remains underfit, and that coverage and calibration still require additional validation before model-comparison claims.

Table~\ref{tab:acoustic-preflight} confirms that the dependency-free feature rows are present and variable. The public reproduction script also computes source-label acoustic anchors: strong-intensity labels have higher RMS than normal labels, and raised-pitch plan labels have higher rough pitch than lowered-pitch labels. These statistics are not a substitute for a perceptual study or planner-ranking evidence; they only check that future acoustic-control experiments have a usable input channel and that selected source labels align with waveform-derived measurements.

\section{Limitations}
The current results are intentionally bounded. The source-label diagnostic is not a claim of broad speech-benchmark superiority, and the model checks still suffer from incomplete valid prediction coverage. The paper also does not evaluate listener judgments, generated-speech user ratings, or real-world release behavior. Those questions require a separate study design with recruited raters and stronger generation provenance. The contribution here is the process layer: before asking downstream user questions, we ask whether the model grounds its reasoning in the evidence it claims to use.

\section{Reproducibility}
The public repository includes the paper source, the 100-case VoxReasonBench source-label split, planner prompt files, gold planner outputs, compact result files, result-reproduction code, scoring code, and tests. Validate the benchmark data, rebuild prompts, score gold test predictions, and recreate the derived result files with:
\begin{verbatim}
python3 scripts/validate_benchmark_data.py
python3 scripts/build_benchmark_prompts.py
python3 scripts/score_predictions.py data/benchmark/source_label/test_gold_predictions.jsonl --split test
python3 scripts/reproduce_results.py
python3 -m pytest
\end{verbatim}
No site-specific job launch scripts, machine-local paths, model weights, raw audio, or raw model completions are included. The released benchmark files contain derived labels, cue records, text prompts, metadata, and expected speaking plans; waveform inputs remain available from the cited RAVDESS source under its license.

\bibliographystyle{plainnat}
\bibliography{references}

\end{document}
